\section{هفته‌ی 16: ارائه‌های پروژه‌ی پایانی و شیوه‌ی خوب مدل‌سازی}

\begin{quote}
\textbf{جمع‌بندی درس} · وقت درس: جمع‌بندی + اولین ارائه‌ها · وقت \lr{lab:}
ارائه‌ها و بازخورد همتایان
\end{quote}

\begin{center}\rule{0.5\linewidth}{0.5pt}\end{center}

\subsection{واژه‌های کلیدی}

واژه‌های مهم این درس با معادل فارسی‌شان (حتی اگر در درس‌های قبل آمده
باشند).

{\def\LTcaptype{none} % do not increment counter
\begin{longtable}[]{@{}
  >{\raggedright\arraybackslash}p{(\linewidth - 4\tabcolsep) * \real{0.3333}}
  >{\raggedright\arraybackslash}p{(\linewidth - 4\tabcolsep) * \real{0.3333}}
  >{\raggedright\arraybackslash}p{(\linewidth - 4\tabcolsep) * \real{0.3333}}@{}}
\toprule\noalign{}
\begin{minipage}[b]{\linewidth}\raggedright
\lr{English}
\end{minipage} & \begin{minipage}[b]{\linewidth}\raggedright
فارسی
\end{minipage} & \begin{minipage}[b]{\linewidth}\raggedright
توضیح کوتاه
\end{minipage} \\
\midrule\noalign{}
\endhead
\bottomrule\noalign{}
\endlastfoot
\lr{peer feedback} & بازخورد همتا & نقد سازنده‌ی کار هم‌کلاسی‌ها \\
\lr{fit} / \lr{fitting} & برازش & تنظیم پارامترها تا مدل با داده جور
شود \\
\lr{parameter recovery} & بازیابی پارامتر & آیا برازش، پارامترهای معلوم
شبیه‌سازی را برمی‌گرداند \\
\lr{reproducibility} & بازتولیدپذیری & امکان به دست آوردن دوباره‌ی همان
نتایج \\
\lr{simulation} & شبیه‌سازی & تولید داده‌ی مصنوعی با اجرای مدل \\
\lr{estimate} / \lr{estimator} & برآورد / برآوردگر & مقدار به‌دست‌آمده
برای پارامتر / روش به دست آوردن آن \\
\lr{parameter} & پارامتر & کمیت قابل تنظیم مدل که بین افراد یا شرایط فرق
می‌کند \\
\lr{likelihood} & درست‌نمایی & احتمال داده به‌عنوان تابعی از پارامترها \\
\lr{Bayes factor (BF)} & عامل بیز \lr{(Bayes factor)} & نسبت
درست‌نمایی‌های حاشیه‌ای دو مدل \\
\lr{bounds} / \lr{constraints} & کران‌ها / قیدها & محدودیت‌های مجاز برای
مقدار پارامترها \\
\lr{posterior} & توزیع پسین \lr{(posterior)} & باور درباره‌ی پارامتر بعد
از دیدن داده \\
\lr{hierarchical model} & مدل سلسله‌مراتبی (چندسطحی) & پارامترهای فردی از
یک توزیع گروهی می‌آیند \\
\lr{adequacy check} & بررسی کفایت مدل & آیا مدل ویژگی‌های اصلی داده را
بازتولید می‌کند \\
\lr{posterior predictive check (PPC)} & بررسی پیش‌بینی پسین & مقایسه‌ی
داده‌ی شبیه‌سازی‌شده از پسین با داده‌ی واقعی \\
\lr{cross-validation (CV)} & اعتبارسنجی متقابل \lr{(cross-validation)} &
برازش روی بخشی از داده و آزمون روی بخش کنارگذاشته \\
\lr{AIC} / \lr{BIC} & معیار اطلاعاتی آکائیکه / بیزی & برازش به‌علاوه‌ی
جریمه برای تعداد پارامتر \\
\lr{research design} & طرح آزمایش & شرایط، تعداد کوشش‌ها و نمونه‌ها \\
\lr{model recovery} & بازیابی مدل & آیا مقایسه‌ی مدل، مدل مولد را درست
انتخاب می‌کند \\
\lr{drift-diffusion model (DDM)} & مدل \lr{drift-diffusion}
(رانش\lr{--}انتشار) & انباشت پیوسته‌ی شواهد بین دو مرز \\
\lr{preregistration} & پیش‌ثبت \lr{(preregistration)} & ثبت برنامه‌ی تحلیل
پیش از دیدن داده \\
\lr{prior} & توزیع پیشین \lr{(prior)} & باور درباره‌ی پارامتر پیش از دیدن
داده \\
\lr{random seed} & \lr{seed} تصادفی & عدد شروع مولد اعداد تصادفی برای
تکرارپذیری \\
\lr{large language model (LLM)} & مدل بزرگ زبانی & شبکه‌ی عصبی‌ای که کلمه‌ی
بعدی متن را پیش‌بینی می‌کند \\
\lr{accuracy} & دقت & نسبت پاسخ‌های درست \\
\end{longtable}
}

\begin{center}\rule{0.5\linewidth}{0.5pt}\end{center}

\subsection{اهداف یادگیری}

در پایان این هفته، دانشجویان می‌توانند:

\begin{enumerate}
\def\labelenumi{\arabic{enumi}.}
\tightlist
\item
  یک مطالعه‌ی مدل‌سازی کامل (سؤال، مدل، برازش، ارزیابی، بازیابی، تفسیر) را
  روشن و در محدوده‌ی زمانی \textbf{ارائه کنند}.
\item
  مطالعه‌های مدل‌سازی دیگران را با معیارهای ارزیابی درس، به شکلی سازنده
  \textbf{نقد کنند}.
\item
  یک چک‌لیست شیوه‌ی خوب مدل‌سازی (بازتولیدپذیری، شفافیت، استواری) را روی
  گزارش پروژه‌ی خودشان \textbf{اعمال کنند}.
\item
  روش‌های درس را در یک گردش کار منسجم \textbf{یکپارچه کنند} و
  \textbf{توضیح دهند} کِی از کدام استفاده کنند.
\end{enumerate}

\begin{center}\rule{0.5\linewidth}{0.5pt}\end{center}

\subsection{مقدمه}

یک مدل فقط به اندازه‌ی توانایی دیگران در فهمیدن، بررسی کردن و ساختن روی
آن مفید است. در 15 هفته، همه‌ی \lr{lab}ها یک مسیر را دنبال کرده‌اند:
\textbf{صورت‌بندی \ensuremath{\leftarrow} شبیه‌سازی \ensuremath{\leftarrow} برازش \ensuremath{\leftarrow} ارزیابی \ensuremath{\leftarrow} بازیابی \ensuremath{\leftarrow} تفسیر.} این
هفته این مسیر را در پروژه‌ی خودتان نشان می‌دهید، و نیمه‌ی دیگر مدل‌سازی علمی
را تمرین می‌کنید: خواندن مدل دیگران به شکلی انتقادی ولی منصفانه. جلسه با
شیوه‌هایی تمام می‌شود که کار مدل‌سازی را قابل اعتماد و قابل استفاده‌ی دوباره
می‌کنند: کد بازتولیدپذیر، گزارش شفاف، بیان صادقانه‌ی محدودیت‌ها.

\begin{center}\rule{0.5\linewidth}{0.5pt}\end{center}

\subsection{جزوه‌ی درس}

\paragraph{به زبان ساده}

هفته‌ی آخر روش جدیدی ندارد. دانشجوها پروژه‌هایشان را ارائه می‌دهند و کل درس
به شکل یک گردش کار خلاصه می‌شود: نظریه را به شکل یک مدل \textbf{فرمال
کن}، آن را \textbf{شبیه‌سازی کن} تا ببینی چه پیش‌بینی می‌کند، آن را روی
داده \textbf{برازش کن}، \textbf{ارزیابی‌اش کن} (آیا برازش خوبی دارد و آیا
از رقیب‌ها بهتر است؟)، بررسی کن که روش در این طرح آزمایش می‌تواند حقیقت را
\textbf{بازیابی کند}، و نتیجه را به زبان روان‌شناختی \textbf{تفسیر کن}،
از جمله اینکه چه نتیجه‌ای را نمی‌شود گرفت. چک‌لیست شیوه‌ی خوب این قدم‌ها را
به چیزهای مشخصی تبدیل می‌کند که در هر گزارش باید بررسی شوند: کد
بازتولیدپذیر، گزارش عدم قطعیت، بررسی کفایت مدل، بازیابی و بیان صادقانه‌ی
محدودیت‌ها.

\textbf{زمان‌بندی پیشنهادی (وقت درس، 60 دقیقه):} جمع‌بندی درس 15 · چک‌لیست
شیوه‌ی خوب 15 · اولین ارائه‌ها 30.

\subsubsection{1. کل درس در یک جدول}

{\def\LTcaptype{none} % do not increment counter
\begin{longtable}[]{@{}
  >{\raggedright\arraybackslash}p{(\linewidth - 4\tabcolsep) * \real{0.3333}}
  >{\raggedright\arraybackslash}p{(\linewidth - 4\tabcolsep) * \real{0.3333}}
  >{\raggedright\arraybackslash}p{(\linewidth - 4\tabcolsep) * \real{0.3333}}@{}}
\toprule\noalign{}
\begin{minipage}[b]{\linewidth}\raggedright
سؤالی که با آن روبه‌رو هستید
\end{minipage} & \begin{minipage}[b]{\linewidth}\raggedright
ابزارها
\end{minipage} & \begin{minipage}[b]{\linewidth}\raggedright
هفته
\end{minipage} \\
\midrule\noalign{}
\endhead
\bottomrule\noalign{}
\endlastfoot
چطور یک نظریه را به مدل تبدیل کنم؟ & پنج جزء؛ شبیه‌سازی؛ بررسی‌های تحلیلی
& 1\lr{--}2 \\
چطور پارامترها را برآورد کنم؟ & کمترین مربعات؛ بیشینه‌ی درست‌نمایی؛
بهینه‌سازی \lr{(grid}، \lr{Nelder\lr{--}Mead}، \lr{L-BFGS-B)}،
\lr{multi-start}، کران‌ها & 3\lr{--}4 \\
برآوردها چقدر مطمئن‌اند؟ & \lr{SE}های \lr{Hessian}، \lr{profile
likelihood}، \lr{bootstrap}؛ بازه‌های پسین & 5، 9 \\
با شرکت‌کننده‌های زیاد چه کنم؟ & برازش‌های فردی؛ خطاهای میانگین‌گیری؛ مدل‌های
سلسله‌مراتبی & 5، 13 \\
آیا مدل کفایت دارد؟ & باقی‌مانده‌ها، \(G^2\)، بررسی‌های پیش‌بینی پسین &
2\lr{--}4، 10 \\
کدام مدل بهتر است؟ & \lr{CV}، \lr{LRT}، \lr{AIC/BIC}، \lr{Bayes factor}
& 6\lr{--}7، 10 \\
آیا طرح من می‌تواند به سؤال جواب دهد؟ & بازیابی پارامتر، بازیابی مدل &
7 \\
چطور برای مدل‌های پیچیده پسین به دست بیاورم؟ & \lr{Grid} \ensuremath{\leftarrow}
\lr{Metropolis\lr{--}Hastings} \ensuremath{\leftarrow} \lr{Gibbs} \ensuremath{\leftarrow} \lr{NUTS (PyMC)}؛ \(\hat R\)،
\lr{ESS}، \lr{divergence}ها & 9\lr{--}13 \\
مدل‌ها در حوزه‌های مختلف چه شکلی‌اند؟ & \lr{SDT}، فراموشی، تمرین، \lr{GCM}،
\lr{DDM/EZ}، \lr{RL}، مشاهده‌گرهای ایده‌آل & 2، 5، 14\lr{--}15 \\
\end{longtable}
}

\paragraph{برای یک مهندس}

با نگاه به عقب، بیشتر درس روی ابزارهایی که قبلاً می‌شناسید منطبق است، و
جدول زیر می‌تواند خلاصه‌ی خودتان برای جواب دادن به سؤال‌های دانشجوها باشد:

{\def\LTcaptype{none} % do not increment counter
\begin{longtable}[]{@{}
  >{\raggedright\arraybackslash}p{(\linewidth - 2\tabcolsep) * \real{0.5000}}
  >{\raggedright\arraybackslash}p{(\linewidth - 2\tabcolsep) * \real{0.5000}}@{}}
\toprule\noalign{}
\begin{minipage}[b]{\linewidth}\raggedright
موضوع درس
\end{minipage} & \begin{minipage}[b]{\linewidth}\raggedright
معادل مهندسی
\end{minipage} \\
\midrule\noalign{}
\endhead
\bottomrule\noalign{}
\endlastfoot
\lr{random walk} / \lr{DDM} (هفته‌های 1، 14) & آشکارسازی ترتیبی
\lr{(SPRT)}، زمان‌های اولین عبور \\
\lr{SDT}، \lr{ROC} (هفته‌ی 2) & آشکارسازی دودویی، \(P_D\) / \(P_{FA}\)،
\lr{ROC} \\
کمترین مربعات، بهینه‌سازها (هفته‌ی 3) & تخمین \lr{LS}، روش‌های گرادیانی،
\lr{Gauss\lr{--}Newton} \\
بیشینه‌ی درست‌نمایی، \(G^2\)، \lr{LRT} (هفته‌های 4، 7) & تخمین \lr{ML}،
آزمون‌های نیکویی برازش، \lr{GLRT} \\
\lr{SE} از \lr{Hessian}، \lr{bootstrap} (هفته‌ی 5) & اطلاعات فیشر /
\lr{CRLB}، ارزیابی کارایی با \lr{Monte Carlo} \\
\lr{CV}، \lr{AIC}، \lr{BIC} (هفته‌های 6\lr{--}7) & انتخاب مرتبه‌ی مدل
\lr{(AIC}، \lr{MDL)}، داده‌ی اعتبارسنجی \\
بازیابی پارامتر و مدل (هفته‌ی 7) & آزمودن تخمین‌گرها روی داده‌ی مصنوعی با
حقیقت معلوم \\
استنباط بیزی، \lr{BF} (هفته‌های 9\lr{--}10) & تخمین \lr{MAP/MMSE}، آشکارسازی
بیزی \\
\lr{MCMC}، \(\hat R\)، \lr{ESS} (هفته‌های 11\lr{--}12) & شبیه‌سازی \lr{Monte
Carlo}، فرایندهای مارکوف، آمار نمونه‌های همبسته \\
مدل‌های سلسله‌مراتبی (هفته‌ی 13) & \lr{empirical Bayes}، \lr{shrinkage} از
نوع \lr{James\lr{--}Stein} \\
قاعده‌ی دلتا، ترکیب نشانه‌ها (هفته‌ی 15) & تطبیق \lr{LMS}، \lr{MRC} / ادغام
حسگر در کالمن \\
\end{longtable}
}

چیزی که در مدل‌سازی شناختی واقعاً متفاوت است: «سیستم» یک انسان است، داده
کم و پرنویز است، و خود پارامترها یافته‌ی علمی‌اند؛ پس تفسیر، عدم قطعیت و
بازیابی از کارایی خام مهم‌ترند.

\subsubsection{2. یک چک‌لیست شیوه‌ی خوب مدل‌سازی}

\textbf{پیش از برازش} - {[} {]} سؤال پژوهشی بر حسب پارامترهای مدل یا
مقایسه‌ی مدل‌ها بیان شده است. - {[} {]} مدل‌های نامزد (دست‌کم با یک جایگزین
ساده‌تر) \emph{پیش از} نگاه کردن به داده مشخص شده‌اند، ترجیحاً با پیش‌ثبت
\lr{(preregistration).} - {[} {]} پیشین‌ها (اگر بیزی است) توجیه شده‌اند و
با شبیه‌سازی پیش‌بینی پیشین بررسی شده‌اند. - {[} {]} بازیابی پارامتر و مدل
\textbf{در همان طرح واقعی} بررسی شده است.

\textbf{برازش و ارزیابی} - {[} {]} درست‌نمایی با نوع داده جور است
(دوجمله‌ای برای شمارش‌ها، ترتیبی برای یادگیری، \ensuremath{\ldots}). - {[} {]}
بهینه‌سازی از \lr{multi-start} و کران‌های معقول استفاده می‌کند و برازش‌ها
به‌صورت تصویری بررسی شده‌اند؛ برای \lr{MCMC}، \(\hat R\)، \lr{ESS} و
\lr{divergence}ها گزارش شده‌اند. - {[} {]} برازش مطلق (باقی‌مانده‌ها،
\(G^2\) یا \lr{PPC)} \emph{و} برازش نسبی \lr{(CV/IC/BF)} گزارش شده‌اند. -
{[} {]} برای هر پارامتری که تفسیر می‌شود عدم قطعیت گزارش شده است.

\textbf{گزارش و بازتولیدپذیری} - {[} {]} همه‌ی کدها و (در صورت اجازه)
داده‌ها به اشتراک گذاشته شده‌اند؛ \lr{seed}های تصادفی ثابت شده‌اند. - {[}
{]} نسخه‌های نرم‌افزار ثبت شده‌اند (مثلاً
\lr{\nolinkurl{conda env export > environment.yml}}). - {[} {]} یک دستور
یا \lr{notebook} همه‌ی عددها و نمودارهای گزارش را دوباره تولید می‌کند. -
{[} {]} محدودیت‌ها بیان شده‌اند: طرح چه چیزی را نمی‌تواند تمیز دهد، کدام
پارامترها قابل بازیابی نیستند، کدام فرض‌ها آزموده نشده‌اند.

\textbf{ثبت محیط محاسباتی} (در هر \lr{notebook} پروژه بگذارید):

\begin{quote}
\textbf{\lr{Code} · کد:} \lr{\nolinkurl{code/week16.py}}, \lr{cell} 1 ·
سلول 1
\end{quote}

\subsubsection{3. قالب ارائه}

\begin{itemize}
\tightlist
\item
  \textbf{مدت:} 8 دقیقه (انفرادی) یا 10 دقیقه (دونفره)، به‌علاوه‌ی 4 دقیقه
  پرسش. زمان دقیق رعایت می‌شود.
\item
  \textbf{ساختار پیشنهادی} (تقریباً هر مورد یک اسلاید):

  \begin{enumerate}
  \def\labelenumi{\arabic{enumi}.}
  \tightlist
  \item
    سؤال و اینکه چرا مهم است.
  \item
    مدل(ها): معادله‌ها و معنای روان‌شناختی هر پارامتر.
  \item
    داده و طرح.
  \item
    روش برازش و یک نمودار برازش (داده + پیش‌بینی‌ها).
  \item
    ارزیابی: مقایسه‌ی مدل \emph{و} یک بررسی کفایت.
  \item
    بازیابی: آیا طرح می‌تواند به سؤال جواب دهد؟
  \item
    تفسیر و محدودیت‌ها.
  \end{enumerate}
\item
  \textbf{تیم‌های دونفره:} هر دو عضو ارائه می‌دهند، و بخش بیزی/\lr{MCMC}
  یا سلسله‌مراتبی باید نشان داده شود.
\item
  \textbf{برنامه:} یک هفته زودتر اعلام می‌شود؛ ترتیب با قرعه تعیین می‌شود.
  جلسه‌ها در وقت درس و \lr{lab} برگزار می‌شوند (حدود 15\lr{--}18 ارائه در 3
  ساعت، بسته به تعداد ثبت‌نام).
\end{itemize}

\subsubsection{4. دادن و گرفتن بازخورد}

هر دانشجو برای سه ارائه‌ی تعیین‌شده یک \textbf{فرم بازخورد همتا} پر می‌کند
(پایین). بازخورد برای ارائه‌دهنده ناشناس و برای مدرس قابل مشاهده است و در
نمره‌ی مشارکت \lr{(lab check-off)} حساب می‌شود. بازخورد خوب \textbf{مشخص}
است («نمودار بازیابی نشان می‌دهد \(\beta\) \lr{identifiable} نیست، پس
تفاوت گروه‌ها در \(\beta\) نباید تفسیر شود») نه کلی («ارائه‌ی خوبی بود»)،
و \textbf{سازنده} است (یک بهبود مشخص پیشنهاد می‌کند).

\textbf{فرم بازخورد همتا}

{\def\LTcaptype{none} % do not increment counter
\begin{longtable}[]{@{}lll@{}}
\toprule\noalign{}
مورد & امتیاز (1\lr{--}4) & توضیح \\
\midrule\noalign{}
\endhead
\bottomrule\noalign{}
\endlastfoot
سؤال روشن بود و مدل با آن جور بود & & \\
پارامترهای مدل با زبان روان‌شناختی توضیح داده شدند & & \\
برازش نشان داده و ارزیابی شد (کفایت + مقایسه) & & \\
به بازیابی یا \lr{identifiability} پرداخته شد & & \\
نتیجه‌گیری‌ها از نتایج بیرون آمدند و محدودیت‌ها بیان شدند & & \\
\textbf{یک چیزی که به‌خصوص خوب انجام شده بود} & & \\
\textbf{یک پیشنهاد مشخص برای گزارش کتبی} & & \\
\end{longtable}
}

\paragraph{برای مدرس}

\begin{itemize}
\tightlist
\item
  \textbf{با \lr{rubric} نمره بدهید و با چک‌لیست نظر بدهید.} در طول
  ارائه‌ها چک‌لیست شیوه‌ی خوب را کنارتان نگه دارید و برای هر تیم یک نقطه‌ی
  قوت و یک پیشنهاد مشخص یادداشت کنید؛ این کار بازخورد کتبی را سریع و
  منصفانه می‌کند.
\item
  \textbf{سؤال‌هایی برای ارائه‌دهندگان} (منطبق بر درس): کدام پارامتر به
  سؤال شما جواب می‌دهد؟ از کجا می‌دانید قابل بازیابی است؟ آیا مدل به‌طور
  مطلق برازش خوبی دارد، نه فقط بهتر از رقیب؟ طرح آزمایش شما چه چیزی را
  \emph{نمی‌تواند} به ما بگوید؟
\item
  \textbf{از نقد حل‌شده} (مطالعه‌ی خیالی \lr{Q-learning)} به‌عنوان الگوی
  بازخورد دادن دانشجوها استفاده کنید: مشخص، وصل به یکی از هفته‌های درس، و
  سازنده.
\item
  \textbf{درس را با جدول خلاصه‌ی یک‌صفحه‌ای ببندید.} از دانشجوها بپرسید
  کدام سطر برایشان سخت‌تر بود و برای مرور امتحان آن‌ها را به هفته‌ی مربوط
  ارجاع دهید.
\end{itemize}

\begin{center}\rule{0.5\linewidth}{0.5pt}\end{center}

\subsection{مثال حل‌شده: نقد یک مطالعه‌ی مدل‌سازی}

\textbf{خلاصه‌ی یک مطالعه‌ی (خیالی).} \emph{«یک مدل \lr{Q-learning} را روی
30 بیمار و 30 فرد کنترل برازش کردیم که 60 \lr{trial} از یک تکلیف
\lr{bandit} دوبازویی با احتمال‌های پاداش ثابت انجام دادند. هر شرکت‌کننده
با \lr{MLE} برازش شد. بیماران نرخ یادگیری به‌طور معنادار پایین‌تری داشتند
(\(t(58) = 2.4\)، \(p = .02\)). نتیجه می‌گیریم که این اختلال یادگیری از
پاداش را مختل می‌کند.»}

\textbf{نقد با استفاده از چک‌لیست.}

\begin{enumerate}
\def\labelenumi{\arabic{enumi}.}
\tightlist
\item
  \textbf{بازیابی در این طرح (هفته‌ی 7، \lr{lab} هفته‌ی 15).} با 60
  \lr{trial} و احتمال‌های پاداش \emph{ثابت}، نرخ یادگیری ضعیف شناسایی
  می‌شود، و \(\alpha\) و \(\beta\) با هم \lr{trade-off} دارند. بدون
  مطالعه‌ی بازیابی، تفاوت در \(\hat\alpha\) ممکن است بازتاب تفاوت در نویز
  انتخاب (\(\beta\)) یا فقط نویز برآورد باشد. \emph{پیشنهاد:} بازیابی
  پارامتر را در همین طرح گزارش کنید و وارونگی اضافه کنید.
\item
  \textbf{مقایسه‌ی مدل (هفته‌ی 15، بخش \lr{D).}} آیا \lr{Q-learning} در هر
  دو گروه از راهبردهای ساده‌تر \lr{(WSLS}، پافشاری) بهتر بود؟ اگر بیماران
  با راهبرد دیگری بهتر توصیف شوند، مقایسه‌ی \(\alpha\) آن‌ها با گروه کنترل
  معنادار نیست. \emph{پیشنهاد:} \lr{AIC/BIC} یا \lr{CV} را برای هر گروه
  همراه با بازیابی مدل گزارش کنید.
\item
  \textbf{آزمون \lr{t} مرحله‌ی دوم روی برآوردهای نقطه‌ای (هفته‌های 5، 13).}
  عدم قطعیت برآوردهای فردی را نادیده می‌گیرد، که با 60 \lr{trial} بزرگ
  است. \emph{پیشنهاد:} یک مدل سلسله‌مراتبی با اثر گروه روی \(\alpha\)، و
  گزارش پسین آن.
\item
  \textbf{کفایت (هفته‌های 4، 10).} هیچ بررسی پیش‌بینی پسین یا نمودار منحنی
  یادگیری نشان داده نشده. آیا مدل منحنی یادگیری هر گروه را بازتولید
  می‌کند؟
\item
  \textbf{تفسیر.} «یادگیری از پاداش را مختل می‌کند» یک ادعای فرایندی قوی
  است. \(\hat\alpha\) پایین‌تر با کاوش بیشتر، انگیزه‌ی کمتر (\(\beta\)
  پایین‌تر) یا پیشین‌های متفاوت روی ارزش‌های اولیه هم سازگار است.
  \emph{پیشنهاد:} نتیجه‌گیری را تعدیل کنید، یا طرحی اضافه کنید که این
  تبیین‌ها را از هم جدا کند.
\end{enumerate}

\textbf{در مجموع:} یک مطالعه‌ی اول معقول که نتیجه‌گیری‌اش هنوز پشتوانه
ندارد. سه افزوده (بازیابی، مقایسه‌ی مدل، برآورد سلسله‌مراتبی) آن را
قانع‌کننده می‌کند. توجه کنید که هر نکته مستقیماً به یکی از هفته‌های این درس
برمی‌گردد.

\begin{center}\rule{0.5\linewidth}{0.5pt}\end{center}

\subsection{\texorpdfstring{\lr{lab} (2 ساعت): ارائه‌ها و بازخورد
همتایان}{lab (2 ساعت): ارائه‌ها و بازخورد همتایان}}

\textbf{اهداف:} پروژه‌تان را ارائه دهید؛ برای سه ارائه بازخورد کتبی
بدهید؛ بازخوردهایی را که می‌گیرید برای گزارش پایانی یادداشت کنید.

\textbf{هماهنگی‌ها:} - اسلایدها \lr{(PDF)} را \textbf{پیش از} شروع جلسه
در صفحه‌ی درس بارگذاری کنید. - لپ‌تاپتان را با \lr{notebook} پروژه
بیاورید، برای پرسش‌هایی درباره‌ی کد. - فرم‌های بازخورد همتا را در طول جلسه
پر کنید (کاغذی یا فرم آنلاین) و تا پایان روز تحویل دهید.

\textbf{بعد از ارائه‌ها:} مدرس به هر تیم بازخورد شفاهی کوتاهی می‌دهد.
بازخورد کتبی (مدرس + همتایان) ظرف 3 روز برگردانده می‌شود تا در گزارش
پایانی استفاده شود.

\begin{center}\rule{0.5\linewidth}{0.5pt}\end{center}

\subsection{\texorpdfstring{ارتباط با هوش مصنوعی و
\lr{LLM}ها}{ارتباط با هوش مصنوعی و LLMها}}

\textbf{شیوه‌ی خوب وقتی هوش مصنوعی بخشی از کار مدل‌سازی شماست.} ابزارهای
هوش مصنوعی امروز با سه نقش وارد مدل‌سازی شناختی می‌شوند، و هر نقش به
عادت‌های بررسی این درس نیاز دارد: 1. \textbf{به‌عنوان دستیار کدنویسی.} کد
تولیدشده می‌تواند درست به نظر برسد و غلط باشد. آن را به روش مدل‌سازی
بیازمایید: داده را با پارامترهای معلوم شبیه‌سازی کنید و بررسی کنید که کد
آن‌ها را بازیابی می‌کند (بازیابی پارامتر به‌عنوان \lr{unit test).} 2.
\textbf{به‌عنوان «شرکت‌کننده‌ی شبیه‌سازی‌شده».} بعضی مطالعه‌ها آزمایش‌های
روان‌شناسی را به \lr{LLM}ها می‌دهند و جواب‌ها را با داده‌ی انسانی مقایسه
می‌کنند، یا حتی مدل‌هایی مبتنی بر \lr{LLM} برای رفتار انسان پیشنهاد می‌کنند
(مثلاً \lr{Binz} \lr{\&} \lr{Schulz,} 2023; \lr{Dillion et al.,} 2023 {[}بررسی
شود{]}). چنین ادعاهایی را مثل هر مدل دیگری ارزیابی کنید: تعمیم به
تکلیف‌های تازه، مقایسه با مدل‌های شناختی ساده‌تر، تقلید \lr{(mimicry)} و
تفسیرپذیری (هفته‌ی 15). 3. \textbf{به‌عنوان موضوع مطالعه.} اگر یک \lr{LLM}
را می‌آزمایید، نام و نسخه‌ی مدل، تاریخ، \lr{prompt}های دقیق، دما و تعداد
نمونه‌ها را ثبت کنید. بدون این‌ها مطالعه قابل بازتولید نیست.

کد پایین یک تابع \lr{SDT} ازپیش‌آماده را که «یک دستیار نوشته» می‌آزماید.
معقول به نظر می‌رسد و حتی حالت \(c = 0\) را درست جواب می‌دهد، ولی یک آزمون
بازیابی نشان می‌دهد که علامت \(c\) را غلط حساب می‌کند.

\begin{quote}
\textbf{\lr{Code} · کد:} \lr{\nolinkurl{code/week16.py}}, \lr{cell} 2 ·
سلول 2
\end{quote}

\textbf{خروجی مورد انتظار:}

\Needspace{8\baselineskip}
\begin{latin}
\begin{Verbatim}
true (d', c) = (1.5, 0.3) -> recovered (1.50, -0.30)  FAIL
true (d', c) = (1.0, -0.4) -> recovered (1.00, 0.40)  FAIL
true (d', c) = (2.0, 0.0) -> recovered (2.00, 0.00)  OK
\end{Verbatim}
\end{latin}

یک مشاهده‌گر \lr{liberal}، محافظه‌کار گزارش می‌شد، که نتیجه‌گیری یک مطالعه‌ی
سوگیری را برعکس می‌کرد. آزمونی که فقط \(c = 0\) را بسنجد این را پیدا
نمی‌کرد.

\textbf{امتحانش کنید (بدون نیاز به \lr{API).}} - \ensuremath{\star} تابع را اصلاح کنید و
آزمون‌ها را دوباره اجرا کنید. - \ensuremath{\star}\ensuremath{\star} یک آزمون بازیابی مشابه برای تابع
\lr{EZ-diffusion} هفته‌ی 14 بنویسید، با استفاده از شبیه‌ساز همان \lr{lab.}
- \ensuremath{\star}\ensuremath{\star} برای پروژه‌ی خودتان: یک پاراگراف به بیانیه‌ی بازتولیدپذیری اضافه کنید
که توضیح دهد ابزارهای هوش مصنوعی چطور استفاده شدند و خروجی‌شان چطور بررسی
شد (\lr{\nolinkurl{assessments/project_brief.md}}، بخش 6 را ببینید).

\emph{مطالعه‌ی اختیاری:} \lr{Binz, M.,} \lr{\&} \lr{Schulz, E.} (2023).
\lr{Using cognitive psychology to understand GPT-3.} \emph{\lr{PNAS,}
120}(6), \lr{e2218523120. Dillion, D., Tandon, N., Gu, Y.,} \lr{\&} \lr{Gray,
K.} (2023). \lr{Can AI language models replace human participants?}
\emph{\lr{Trends in Cognitive Sciences,} 27}, 597\lr{--}600 {[}بررسی شود{]}.

\begin{center}\rule{0.5\linewidth}{0.5pt}\end{center}

\subsection{اشتباهات رایج و برداشت‌های نادرست (در ارائه‌ها و گزارش‌های
پروژه)}

\begin{enumerate}
\def\labelenumi{\arabic{enumi}.}
\tightlist
\item
  \textbf{پرداختن زیاد به روش و کم به معنا.} هر پارامتر و هر نتیجه به یک
  تفسیر روان‌شناختی نیاز دارد.
\item
  \textbf{نشان دادن یک آماره‌ی برازش بدون نمودار برازش.} همیشه داده را در
  برابر پیش‌بینی‌ها رسم کنید.
\item
  \textbf{مقایسه‌ی مدل بدون بررسی کفایت،} یا برعکس.
\item
  \textbf{تفسیر پارامترهایی که قابل بازیابی بودنشان نشان داده نشده.}
\item
  \textbf{ادعای بیش از حد:} «مدل درست است»، «ذهن بیزی است». بگویید شواهد
  از چه چیزی پشتیبانی می‌کند و از چه چیزی نه.
\item
  \textbf{نتایج غیرقابل بازتولید:} \lr{seed}های جاافتاده، نسخه‌های
  ثبت‌نشده‌ی بسته‌ها، سلول‌های \lr{notebook} که بی‌ترتیب اجرا شده‌اند.
\item
  \textbf{رد کردن از زمان.} ارائه را تمرین کنید؛ در هر اسلاید یک ایده.
\end{enumerate}

\begin{center}\rule{0.5\linewidth}{0.5pt}\end{center}

\subsection{خودآزمایی (پرسش‌های مرور یکپارچه برای امتحان پایانی)}

\textbf{1.} همکاری داده‌ی دقت و \lr{RT} از یک دست‌کاری سرعت/دقت دارد و
می‌خواهد بداند دست‌کاری روی احتیاط اثر می‌گذارد یا کیفیت شواهد. یک برنامه‌ی
کامل مدل‌سازی طرح کنید و مدل، روش برازش، ارزیابی و بررسی بازیابی را نام
ببرید. \textbf{پاسخ} مدل: \lr{drift-diffusion (EZ-diffusion} برای دور
اول، یا \lr{DDM} کامل با \lr{MLE} / بیزی سلسله‌مراتبی). برای هر شرکت‌کننده
و شرط برازش کنید. مدل‌هایی را مقایسه کنید که در آن‌ها دست‌کاری فقط روی
\(a\)، فقط روی \(v\)، یا روی هر دو اثر می‌گذارد \lr{(AIC/BIC} یا یک مدل
سلسله‌مراتبی با اثر شرط؛ \lr{PPC} روی \lr{quantile}های \lr{RT).} بازیابی:
با هر مدل در طرح واقعی داده شبیه‌سازی کنید و بررسی کنید که پارامترها و
نتیجه‌ی «کدام پارامتر تغییر کرد» بازیابی می‌شوند. اثر شرط روی \(a\) و
\(v\) را با عدم قطعیتش تفسیر کنید.

\textbf{2.} کِی برای مقایسه‌ی دو مدل \lr{Bayes factor} را به
\lr{cross-validation} ترجیح می‌دهید، و برعکس؟ \textbf{پاسخ} \lr{Bayes
factor:} وقتی سؤال این است که با توجه به داده کدام مدل محتمل‌تر است
(شواهد، از جمله پشتیبانی از فرض صفر)، وقتی پیشین‌های معنادار در دسترس‌اند،
و همراه با تحلیل حساسیت به پیشین. \lr{CV:} وقتی هدف دقت پیش‌بینی برای
داده‌ی تازه از یک نوع مشخص است، توجیه پیشین‌ها سخت است، یا \lr{BF} بیش از
حد به پیشین‌ها حساس است.

\textbf{3.} توضیح دهید چطور یک پارامتر می‌تواند خیلی دقیق برآورد شود
\lr{(CI} یا پسین باریک) و با این حال بی‌معنی باشد. \textbf{پاسخ} دقت فقط
انحنای درست‌نمایی را تحت مدل فرض‌شده نشان می‌دهد. اگر مدل بد مشخص شده باشد
(در \lr{PPC/}\(G^2\) رد شود)، اگر پارامتر با پارامتر دیگری معاوضه شود
طوری که معنایش به یک قید دلبخواه بستگی داشته باشد، یا اگر بهینه‌ساز یا
نمونه‌گیر شکست خورده باشد (کمینه‌ی محلی، کران، عدم همگرایی)، برآورد دقیق
سازه‌ی مورد نظر را نمی‌سنجد.

\textbf{4.} چرا مدل‌های سلسله‌مراتبی اغلب برآوردهای فردی بهتری از برازش‌های
فردی جداگانه می‌دهند؟ \textbf{پاسخ} از طریق توزیع گروه بین شرکت‌کننده‌ها
اطلاعات به اشتراک می‌گذارند و برآوردهای فردی پرنویز را به نسبت عدم
قطعیتشان به سمت میانگین گروه \lr{shrink} می‌کنند. این واریانس برآورد را
بیشتر از سوگیری‌ای که اضافه می‌کند کم می‌کند، پس برآوردها به‌طور میانگین به
حقیقت نزدیک‌ترند.

\textbf{5.} سه راه مختلفی را که این درس از شبیه‌سازی استفاده کرد نام
ببرید و بگویید چرا هرکدام مهم است. \textbf{پاسخ} (1) کاوش پیش‌بینی‌های مدل
پیش از برازش (چه چیزی را می‌تواند و نمی‌تواند تولید کند). (2) بازیابی
پارامتر و مدل (آیا طرح و روش می‌توانند به سؤال جواب دهند). (3)
\lr{bootstrap} و بررسی‌های پیش‌بینی (عدم قطعیت و کفایت). همچنین: \lr{Monte
Carlo} و \lr{MCMC} برای محاسبه‌ی پسین‌ها.

\textbf{6.} یک مدل بین چهار نامزد بهترین \lr{AIC} را دارد ولی بررسی
پیش‌بینی پسین یک ناهمخوانی نظام‌مند در تأخیرهای طولانی نشان می‌دهد. چه
گزارش می‌کنید و چه نتیجه می‌گیرید؟ \textbf{پاسخ} هر دو را گزارش کنید. مدل
برای پیش‌بینی بهترینِ نامزدهاست، ولی به‌طور مطلق کفایت ندارد: هیچ‌کدام از
نامزدها رفتار تأخیرهای طولانی را نمی‌گیرد. نتیجه‌گیری‌ها باید به همین نسبت
محدود شوند، و قدم بعدی بازنگری مدل با هدف همان ویژگی است.

\begin{center}\rule{0.5\linewidth}{0.5pt}\end{center}

\subsection{تکلیف}

\subsubsection{پروژه‌ی پایانی: گزارش کتبی و کد}

\begin{itemize}
\tightlist
\item
  \textbf{مهلت:} یک هفته بعد از این جلسه، و در هر حال پیش از امتحان
  پایانی (تاریخ دقیق در صفحه‌ی درس).
\item
  \textbf{تحویل:} گزارش \lr{(PDF}، حداکثر 10 صفحه برای انفرادی و 14 برای
  دونفره، بدون منابع و پیوست) به‌علاوه‌ی یک آرشیو کد یا لینک
  \lr{repository} که همه‌ی نتایج را بازتولید کند.
\item
  \textbf{نمره‌دهی:} شرح پروژه و \lr{rubric} را ببینید
  (\lr{\nolinkurl{assessments/project_brief.md}}). گزارش و کد 3 نمره از
  6 نمره‌ی پروژه را دارند؛ پیشنهاد (1) و ارائه (2) قبلاً نمره گرفته‌اند.
\item
  \textbf{از بازخورد امروز استفاده کنید:} \lr{rubric} به پاسخ‌های قابل
  مشاهده به بازخورد ارائه نمره می‌دهد.
\end{itemize}

\subsubsection{امتحان پایانی}

کتبی، جامع (با تأکید بر فصل‌های 4\lr{--}6)، با برگه‌ی فرمول، بدون لپ‌تاپ؛ 3 نمره
از 20. برای قالب و یک نمونه سؤال
\lr{\nolinkurl{assessments/final_exam.md}} را ببینید.

\begin{center}\rule{0.5\linewidth}{0.5pt}\end{center}

\subsection{مطالعه‌ی بیشتر}

\begin{itemize}
\tightlist
\item
  \textbf{\lr{Farrell} \lr{\&} \lr{Lewandowsky} (2018)}: فصل استفاده از مدل‌ها
  در روان‌شناسی (نقش مدل‌ها در توسعه‌ی نظریه و دام‌های رایج) {[}شماره‌ی فصل
  بررسی شود{]}.
\item
  \textbf{\lr{Lee} \lr{\&} \lr{Wagenmakers} (2014)}: مطالعه‌های موردی بخش
  \lr{IV} را برای ایده‌هایی درباره‌ی گسترش پروژه‌تان مرور کنید.
\item
  \textbf{\lr{Martin, Kumar} \lr{\&} \lr{Lao} (2021)، فصل 9.} گردش‌های کار
  بیزی از ابتدا تا انتها، الگویی خوب برای ساختار دادن به یک گزارش.
\item
  \emph{اختیاری، خارج از فهرست درس:} \lr{Wilson, R. C.,} \lr{\&} \lr{Collins,
  A. G. E.} (2019). \lr{Ten simple rules for the computational modeling
  of behavioral data.} \emph{\lr{eLife,} 8}, \lr{e49547.}
\end{itemize}

\textbf{از همراهی‌تان در این درس سپاسگزاریم؛ به مدل‌سازی ادامه دهید!}
